%%%%%%%%%%%%%%%%%%%%%%%%%%%%%%%%%%%%%%%%%%%%%%%%%%%%%%%%%%%%%%%%%%%%%%%%%%%
%%  introduction_v10.tex  –  Введение (ненумерованная глава)
%%  v10 – Обновлено в соответствии с финальными слайдами презентации.
%%       Новое название диссертации: «Разработка состязательных устойчивых
%%       моделей в гибридных системах обнаружения и предотвращения вторжений»
%%       Научная новизна переписана как концепции/теории (не модели).
%%       Объект/предмет обновлены для гибридных IDPS.
%%       Цитирование начинается с [1].
%%%%%%%%%%%%%%%%%%%%%%%%%%%%%%%%%%%%%%%%%%%%%%%%%%%%%%%%%%%%%%%%%%%%%%%%%%%

\chapter*{ВВЕДЕНИЕ}
\addcontentsline{toc}{chapter}{Введение}
\markboth{ВВЕДЕНИЕ}{ВВЕДЕНИЕ}

%%---------------------------------------------------------------------
\section*{Актуальность темы исследования}
%%---------------------------------------------------------------------

Графовые нейронные сети (ГНС)~\cite{Kipf2017,Velickovic2018,Hamilton2017,Wu2020gnn_survey,Luo2021} стали основной парадигмой глубокого обучения для извлечения предсказательных представлений из реляционных данных. Модели ГНС внедряются в критически важные системы: здравоохранение (моделирование молекул как графов атомов и связей, графы взаимодействия пациентов для рекомендаций по лечению), финансы (анализ графов транзакций между счетами для обнаружения мошенничества и отмывания денег), кибербезопасность (анализ сетевого трафика как графов связанных устройств для обнаружения вторжений), промышленные системы (графы слияния датчиков для мониторинга SCADA и АСУ ТП) и автономные системы (графы слияния датчиков для восприятия окружения). Однако устойчивость моделей ГНС отстаёт от их точности на 25--65 процентных пунктов в условиях атак, а именно механизм передачи сообщений, придающий этим моделям их выразительную способность, делает их уязвимыми к возмущениям~\cite{Zuegner2018,Bojchevski2019,Dai2018,Goodfellow2014gan,Szegedy2014}.

Существующие модели ГНС страдают от четырёх взаимосвязанных проблем: (1)~легко обманываются малыми целенаправленными изменениями входных данных или связей узлов; (2)~не способны обрабатывать новые или меняющиеся паттерны без полного переобучения; (3)~слишком медленны и вычислительно затратны для использования в реальном времени на больших графах; и (4)~не способны передавать знания между организациями без утечки конфиденциальных данных. Эти проблемы динамичны по своей природе --- они изменяются во времени, однако имеют общий архитектурный корень.

Чтобы понять, почему возникает эта чувствительность, необходимо осознать двойственную природу входов моделей ГНС. В отличие от обычных нейронных сетей, принимающих единый тензор данных, модель ГНС одновременно получает два входа: числовые \emph{входные данные} ($\mathbf{X}\in\mathbb{R}^{n\times d}$, матрица признаков, кодирующая атрибуты сущностей) и реляционную \emph{входную структуру} ($\mathbf{A}$, матрица смежности, кодирующая попарные связи). Сеть формирует свой выход путём распространения данных вдоль структуры через итеративную агрегацию окрестностей. Возмущение любого из входов --- данных или структуры --- может нарушить предсказания. Именно эта двойственная природа входов отличает модели ГНС от обычных нейронных сетей и является источником проблем робастности, решаемых в настоящей диссертации.

\begin{figure}[H]
\centering
\includegraphics[width=0.89\textwidth]{figures/fig1a_gnn.pdf}
\caption{Двойственная природа входов моделей ГНС. Модель одновременно получает числовые данные (матрица признаков $\mathbf{X}$) и реляционную структуру (матрица смежности $\mathbf{A}$) в качестве входов. Возмущение любого из входов может нарушить обученное представление, что мотивирует разработку системы робастности, реализуемой в гибридных системах обнаружения и предотвращения вторжений, являющихся предметом настоящей диссертации.}
\label{fig:dual_input}
\end{figure}

Большинство моделей ГНС работают в дискретном времени, что усугубляет эти уязвимости. Модели дискретного времени принудительно привязывают нерегулярные события графа к фиксированной тактовой сетке, теряя информацию и гарантии робастности. Формулировки непрерывного времени --- например, через нейронные обыкновенные дифференциальные уравнения~\cite{Chen2018,Chamberlain2021,Rusch2022} --- гладко следуют за временем событий, захватывают множественные временные масштабы в одной модели и обнаруживают границу Липшица, дающую сертифицированную робастность --- преимущество, которое ни одна дискретная модель не может обеспечить. Поскольку представления эволюционируют согласно дифференциальному уравнению с ограниченной константой Липшица, малые возмущения входа могут вызвать лишь ограниченные, гладкие изменения выхода --- гладкость управляющего потока математически ограничивает степень изменения предсказания злоумышленником. Неравенство Грёнуолла формализует это: расстояние между двумя траекториями ОДУ растёт не более чем экспоненциально по $L_g T$, давая замкнутый сертификат робастности:
\begin{equation}
\|\mathbf{h}_1(T)-\mathbf{h}_2(T)\|\le e^{L_g T}\|\mathbf{x}_1-\mathbf{x}_2\|,
\label{eq:gronwall}
\end{equation}
где $\mathbf{h}_i(T)$ --- скрытое состояние траектории $i$ в момент $T$, $\mathbf{x}_i$ --- начальный вход (признаки и структура), $L_g$ --- константа Липшица векторного поля ОДУ, $T$ --- горизонт интегрирования, $e^{L_gT}$ --- максимальный коэффициент усиления, а $\|\cdot\|$ обозначает евклидову норму. Модели ГНС дискретного времени не имеют такого непрерывного механизма, поэтому возмущения непредсказуемо накапливаются от слоя к слою. Однако модели непрерывного времени получают ограниченное внимание в литературе и текущей практике, поэтому настоящая диссертация сосредоточена на их разработке через модели на основе CT-GNN.

Актуальность этих уязвимостей подтверждается тремя сходящимися трендами. Во-первых, масштаб проблемы: глобальный ущерб от кибератак достигает \$10,5 трлн в год к 2025 году, и ИИ-инструменты центров управления безопасностью (SOC) необходимы, но уязвимы. Во-вторых, регулирование: NIST AI RMF, EU AI Act и сроки перехода на постквантовую криптографию требуют сертифицированной устойчивости, а не только эмпирического тестирования. В-третьих, теоретический пробел: нет единой математической основы, одновременно обеспечивающей устойчивость, приватность, эффективность и точность на динамических графовых структурах. Настоящая диссертация восполняет этот пробел: 7 методов с формальными гарантиями по 3 условиям и 4 требованиям.

Документальным подтверждением служат недавние рецензируемые исследования, фиксирующие серьёзные отказы моделей ГНС в четырёх критически важных областях. В здравоохранении градиентная атака на чувствительность рёбер на графовые изоморфные сети вызвала падение AUC приблизительно на 25\% на бенчмарках MoleculeNet (Tox21, ClinTox, SIDER), при этом токсичные соединения были классифицированы как безопасные~\cite{EdgeSensitivity2025}. В финансах атака MonTi с внедрением группы почти утроила ошибку классификации ГНС-детекторов мошенничества (CARE-GNN, PC-GNN, GAGA) с 25,7\% до 68,6\%~\cite{MonTi2025}. В промышленном IoT отравление данных федеративных детекторов вторжений вызвало коллапс точности на 64 процентных пункта с 94,2\% до 30,0\% при не-IID условиях~\cite{FerragPoisoning2023}. В кибербезопасности атака PROVNINJA снизила обнаружение ГНС-систем на основе провенанса до 59\%, сделав APT-атакующих невидимыми~\cite{PROVNINJA2023}. Эти задокументированные отказы подтверждают, что робастность моделей ГНС отстаёт от точности предсказаний на 25--65 процентных пунктов во всех критически важных областях, и этот разрыв мотивирует всю диссертацию.

Распределённая среда --- это среда, в которой несколько организаций (например, больницы, банки, центры обработки данных) владеют частью одних и тех же графовых данных и не могут ими делиться, поэтому общая модель должна обучаться на каждой части через федеративное обучение~\cite{McMahan2017,Blanchard2017,Kairouz2021,Konecny2016,Mothukuri2021} --- без доступа какой-либо стороны к сырым данным другой. Модели ГНС в распределённых средах должны быть более устойчивыми, эффективными, робастными и точными, поскольку они внедряются в критически важные системы быстрее, чем могут быть установлены их гарантии: злоумышленники могут обманывать их малыми возмущениями; распределения смещаются после развёртывания; законы о приватности ограничивают обмен данными~\cite{Dwork2006,Abadi2016,Dwork2014,Mironov2017}; при этом периферийные устройства требуют вывода в реальном времени. Ни один существующий подход не решает все четыре проблемы одновременно --- этот разрыв мотивирует диссертацию.

Реальные графовые данные дрейфуют после развёртывания --- топологии эволюционируют, паттерны атак меняются, появляются новые протоколы~\cite{Gama2014,Lu2019} --- и статические модели ГНС незаметно теряют точность без предупреждения. Это недопустимо в критически важных системах, где устаревшая модель является опасной моделью. Цель состоит в создании динамических моделей ГНС, которые обнаруживают сдвиг распределения онлайн, адаптируют параметры без забывания ранее изученных паттернов и оценивают неопределённость через байесовские апостериорные обновления~\cite{Gal2016,Kendall2017,Lakshminarayanan2017} для маркировки новых входов для экспертного анализа.

Существующие модели ГНС и модели на основе трансформеров~\cite{Vaswani2017,Ying2021} масштабируются квадратично ($\mathcal{O}(L^2)$) с длиной последовательности, поскольку каждый узел должен обращать внимание на каждый другой узел, что делает их слишком медленными и ресурсоёмкими для вывода в реальном времени на больших темпоральных графах с миллионами событий. Литература показывает, что структурированные модели пространства состояний~\cite{Gu2022,Gu2024,Dao2024,Behrouz2024graphmamba} более эффективны, снижая стоимость до $\mathcal{O}(L\log L)$ через свёртки быстрого преобразования Фурье и обеспечивая развёртывание на периферийных устройствах, где модели масштаба трансформеров неприменимы. Поэтому в настоящей диссертации предлагается перенос фокуса внимания и практик с трансформеров на селективные модели пространства состояний для графовых данных.

Описанные выше четыре проблемы не являются независимыми; они образуют логическую цепочку, звенья которой определяют поэтапный охват настоящей диссертации. Динамика графов непрерывного времени обеспечивает основу представлений; федеративные алгоритмы обеспечивают распределённое развёртывание; эффективный вывод на основе пространства состояний позволяет работу в реальном времени; а онлайн-байесовская адаптация гарантирует долгосрочную надёжность. Выход каждой задачи является входом для следующей, и ни один существующий подход не решает все четыре задачи совместно при трёх рабочих условиях и четырёх требованиях, определённых ниже.

%%---------------------------------------------------------------------
%%  Условия и требования
%%---------------------------------------------------------------------

Три рабочих условия и четыре сквозных требования независимы --- они определяют пространство, в котором любая графовая ИИ-система должна функционировать надёжно. В диссертации предлагаются методы, удовлетворяющие всем условиям одновременно.

Три рабочих условия: (1)~состязательная среда, в которой целенаправленный противник возмущает топологию графа, признаки узлов или временной порядок событий; (2)~распределённая среда, в которой графовые данные фрагментированы между организациями, не имеющими возможности централизовать свои данные, и где некоторые участники могут быть ненадёжными или злонамеренными; и (3)~нестационарность, при которой статистические свойства графа эволюционируют во времени.

Четыре сквозных требования: (Т1)~робастность --- предсказания должны оставаться устойчивыми при возмущениях; (Т2)~точность --- должна поддерживаться высокая производительность классификации; (Т3)~эффективность --- события графа должны обрабатываться достаточно быстро для работы в реальном времени; и (Т4)~дифференциальная приватность --- процесс обучения не должен раскрывать конфиденциальные данные.

%%---------------------------------------------------------------------
%%  Таблицы: Условия и требования с назначением моделей
%%---------------------------------------------------------------------

\begin{table}[htbp]
\centering
\caption{Три рабочих условия и методы, разработанные для их решения.}
\label{tab:conditions_methods}
\small
\begin{tabular}{p{2.8cm} p{3.2cm} p{7.5cm}}
\toprule
\textbf{Условие} & \textbf{Методы} & \textbf{Обоснование} \\
\midrule
\textbf{У1: Состязательная среда}
& М1: CT-TGNN\newline М5: Стох.\ трансф.\newline М7: Сертификация
& Нейронные ОДУ непрерывного времени обеспечивают гладкость траектории представления с ограниченной константой Липшица, предоставляя детерминированные сертификаты робастности, которые альтернативы дискретного времени обеспечить не могут. Стохастическое состязательное обучение и теоретико-игровые стратегии равновесия дополняют архитектурную гарантию адаптивной защитой. \\
\addlinespace
\textbf{У2: Нестационарные условия}
& М1: CT-TGNN\newline М3: TripleE-TGNN\newline М4: MambaShield\newline М6: UC-HGP
& CT-TGNN захватывает темпоральную эволюцию через нейронные дифференциальные уравнения; TripleE-TGNN предотвращает катастрофическое забывание через упругую консолидацию весов; MambaShield обеспечивает близкую к линейной эффективность; а UC-HGP предоставляет неопределённость на иерархических гауссовских процессах для обнаружения дрейфа. \\
\addlinespace
\textbf{У3: Распределённая сетевая динамика}
& М2: FedLLM-API\newline М7: Сертификация
& Оптимальный транспорт через итерации Синкхорна корректирует гетерогенность между фрагментами графов, а покоординатная агрегация усечённым средним устойчива к повреждённым участникам, при этом калиброванный гауссовский шум обеспечивает формальный бюджет приватности. \\
\bottomrule
\end{tabular}
\end{table}

\begin{table}[htbp]
\centering
\caption{Четыре сквозных требования и методы, их выполняющие.}
\label{tab:requirements_methods}
\small
\begin{tabular}{p{2.3cm} p{3.5cm} p{7.5cm}}
\toprule
\textbf{Требование} & \textbf{Методы} & \textbf{Обоснование} \\
\midrule
\textbf{Т1: Робастность}
& М1: CT-TGNN\newline М4: MambaShield\newline М7: Сертификация
& CT-TGNN обеспечивает непрерывную динамику с ограниченной константой Липшица; MambaShield устойчив к отравлению на этапе обучения через прогрессивную дистилляцию; система сертификации связывает гарантии через рандомизированное сглаживание и распространение интервальных границ. \\
\addlinespace
\textbf{Т2: Точность}
& М3: TripleE-TGNN\newline М5: Стох.\ трансф.\newline М6: UC-HGP
& Многоуровневые вложения графов захватывают структуру на уровнях сервисов, трассировок и узлов; вариационное байесовское внимание декомпозирует неопределённость; иерархические гауссовские процессы обеспечивают калиброванные доверительные интервалы. \\
\addlinespace
\textbf{Т3: Эффективность}
& М4: MambaShield\newline М2: FedLLM-API\newline М3: TripleE-TGNN
& MambaShield снижает стоимость вывода с квадратичной до близкой к линейной; FedLLM-API сокращает число раундов коммуникации; TripleE-TGNN достигает экономии через структурированную разреженность и обучение с учётом квантования. \\
\addlinespace
\textbf{Т4: Приватность}
& М2: FedLLM-API\newline М1: CT-TGNN\newline М7: Сертификация
& FedLLM-API обеспечивает $(\epsilon,\delta)$-дифференциальную приватность; CT-TGNN поддерживает дифференциально приватное извлечение признаков; система сертификации отслеживает кумулятивный бюджет приватности. \\
\bottomrule
\end{tabular}
\end{table}

Три условия не являются взаимоисключающими: реальные графовые системы, как правило, сталкиваются с двумя или всеми тремя одновременно. На Рис.~\ref{fig:conditions_venn} показано, как семь методов расположены в пространстве, определённом пересекающимися условиями.

\begin{figure}[!t]
\centering
\includegraphics[width=0.89\columnwidth]{figures/fig2a_t_by_f.pdf}
\caption{Расположение семи методов в пространстве условий функционирования (слева) и требований (справа). M4: MambaShield находится на пересечении всех трёх условий; M1: CT-TGNN связывает состязательное и нестационарное условия через сертифицированные непрерывно-временные динамики.}
\label{fig:conditions_venn}
\end{figure}

%%---------------------------------------------------------------------
%%  Архитектурная основа
%%---------------------------------------------------------------------

\begin{figure}[!t]
\centering
\includegraphics[width=0.96\textwidth]{figures/Main_arch_v1a.pdf}
\caption{Архитектурная основа исследовательской программы. Семь методов покрывают все 12 пересечений матрицы условие--требование $3\times4$. Каждый сформулирован на уровне математических операций над графами --- предметно-независим. Все модели разработаны автором и описаны в 18~рецензируемых публикациях.}
\label{fig:dissertation_architecture}
\end{figure}

Система обеспечивает предметную независимость. Каждый метод сформулирован на уровне математических операций над графовыми структурами без привязки к какой-либо конкретной прикладной области. Лишь в последующих главах эти методы отображаются на сущности и ограничения конкретных сред развёртывания, в частности гибридных систем обнаружения и предотвращения вторжений. После определения пространства задач, рабочих условий и архитектурной основы в следующем разделе формулируется точная цель исследования и последовательные задачи для её достижения.


%%---------------------------------------------------------------------
\section*{Цель исследования}
%%---------------------------------------------------------------------

Основной целью настоящей диссертации является повышение робастности моделей гибридных систем обнаружения и предотвращения вторжений (IDPS) к состязательным возмущениям, их устойчивости в нестационарных условиях, эффективности и точности для использования в реальном времени, а также обеспечение приватности в распределённых средах.

Достижение этой цели требует выполнения семи исследовательских задач, расположенных в последовательную цепочку, так что выход каждой задачи обеспечивает вход, необходимый для следующей.

\noindent 1.\quad Разработать модели непрерывного времени (НВ) на основе нейронных обыкновенных дифференциальных уравнений с многомасштабной темпоральной динамикой и сертифицированной робастностью по Липшицу.

\noindent 2.\quad Разработать методы, позволяющие НВ-моделям функционировать в распределённых средах на основе федеративного обучения с византийско-устойчивой агрегацией, дифференциальной приватностью и выравниванием распределений методом оптимального транспорта.

\noindent 3.\quad Разработать методы снижения стоимости компромиссов и повышения эффективности как монолитных, так и распределённых моделей на основе многоуровневых вложений в сочетании с селективными архитектурами пространства состояний (MambaShield), снижающими сложность с квадратичной до близкой к линейной.

\noindent 4.\quad Разработать методы оценки и адаптации моделей ГНС к нестационарным средам --- то есть к дрейфу распределений --- на основе иерархических гауссовских процессов для калиброванной оценки неопределённости и онлайновых байесовских апостериорных обновлений.

\noindent 5.\quad Интегрировать все методы в единую систему, провести валидацию на крупномасштабных бенчмарках и доказать её качество для готовности к развёртыванию с помощью регуляризации согласованности и комплексного тестирования.

\noindent 6.\quad Адаптировать единую систему для задач кибербезопасности, учитывая специфику данной области применения, с помощью предметно-ориентированных адаптеров и уточнения модели угроз.

\noindent 7.\quad Применить единую систему и продемонстрировать её качество в области гибридных систем обнаружения и предотвращения сетевых вторжений (СОПСВ) путём развёртывания в реальном времени на программной платформе RobustIDPS.ai.


%%---------------------------------------------------------------------
\section*{Объект исследования}
%%---------------------------------------------------------------------

Объектом исследования являются состязательно устойчивые модели гибридных систем обнаружения и предотвращения вторжений (IDPS), функционирующие в состязательных, нестационарных и распределённых средах. В конкретных терминах исследование охватывает класс моделей глубокого обучения на основе графовых нейронных сетей, алгоритмов обучения и поддерживающих математических основ, свойства которых определяют, может ли гибридная модель надёжно функционировать при состязательном возмущении топологии графа и признаков узлов, распределительной нестационарности эволюционирующего графа, ограничениях приватности на распределённых фрагментах графа и ограниченных вычислительных бюджетах для обработки темпоральных графов. Это включает нейронные обыкновенные дифференциальные уравнения~\cite{Chen2018,Pontryagin1962} совместно с обучением на основе сопряжённого метода для графовых областей, темпоральные графовые нейронные сети с динамикой непрерывного времени~\cite{Xu2020tgat,Rossi2020,Chamberlain2021}, формулировки оптимального транспорта~\cite{Villani2009,Peyre2019,Cuturi2013,Chizat2020}, структурированные модели пространства состояний~\cite{Gu2022,Dao2024,Behrouz2024graphmamba}, байесовские нейронные сети~\cite{Blundell2015,Lakshminarayanan2017}, теоретико-игровые модели~\cite{Tambe2011,Sinha2018}, а также формулировки робастной оптимизации~\cite{Madry2018,Zhang2019trades,Goodfellow2015}.


%%---------------------------------------------------------------------
\section*{Предмет исследования}
%%---------------------------------------------------------------------

Предметом исследования являются методы обеспечения робастности, эффективности, дифференциальной приватности и точности гибридных моделей, функционирующих в состязательных, нестационарных и распределённых средах. Эти методы включают архитектуры на основе нейронных обыкновенных дифференциальных уравнений для сертифицированной динамики непрерывного времени, алгоритмы федеративной оптимизации с византийско-устойчивой агрегацией и выравниванием на основе оптимального транспорта, многоуровневые темпоральные вложения графов с упругой консолидацией весов, селективные модели пространства состояний для эффективного последовательного вывода, механизмы байесовских апостериорных обновлений для онлайн-адаптации к дрейфу, а также системы теоретико-игровой сертификации для доказуемых гарантий робастности. Математические формулировки и алгоритмические решения сконструированы для переносимости в любую область, где модели гибридных IDPS сталкиваются с состязательными или нестационарными условиями.


%%---------------------------------------------------------------------
\section*{Методы исследования}
%%---------------------------------------------------------------------

Методологическая программа настоящей диссертации организована вокруг восьми формальных методов, каждый из которых выбран потому, что он решает конкретный и измеримый аспект проблемы робастности гибридных моделей IDPS.

Нейронные обыкновенные дифференциальные уравнения моделируют непрерывную эволюцию представлений узлов через параметризованное векторное поле,
\begin{equation}
\frac{d\bh(t)}{dt} = f_\theta\!\bigl(\bh(t),\, t\bigr),
\label{eq:node}
\end{equation}
решаемое численными интеграторами с адаптивным шагом. Их роль --- обеспечить динамику на графовых структурах, чувствительность выхода которой, измеряемая константой Липшица обученного отображения, может быть оценена и представлена как сертификат робастности (центральный элемент методов М1: CT-TGNN и М3: TripleE-TGNN).

Графовые нейронные сети расширяют обученные представления на реляционные данные через агрегацию окрестностей,
\begin{equation}
\bh_v^{(k+1)}
= \phi\!\Bigl(\bh_v^{(k)},\;
  \bigoplus_{u\in\calN(v)}
    \psi\!\bigl(\bh_u^{(k)},\, \mathbf{e}_{uv}\bigr)\Bigr),
\label{eq:gnn_mp}
\end{equation}
и их сочетание с нейронным ОДУ на изменяющейся во времени структуре смежности порождает темпоральные графовые модели, захватывающие как топологические, так и временные зависимости (используется в методах М1, М3 и М5).

Оптимальный транспорт обеспечивает геометрию на распределениях вероятностей через расстояние Вассерштейна,
\begin{equation}
W_p(\mu,\nu)
= \Bigl(\inf_{\gamma\in\Pi(\mu,\nu)}
    \int c(\bx,\by)\,d\gamma(\bx,\by)
  \Bigr)^{\!1/p},
\label{eq:wasserstein}
\end{equation}
чья энтропийная регуляризация через итерации Синкхорна снижает вычислительную стоимость. В данной работе метод применяется для выравнивания гетерогенных распределений фрагментов графов в рамках цикла федеративной оптимизации при сохранении измеримого бюджета дифференциальной приватности (центральный элемент метода М2: FedLLM-API).

Дифференциальная приватность формализует конфиденциальность на уровне записей данных через условие
\begin{equation}
\Prob\bigl[\calM(\calD)\in S\bigr]
\;\le\;
e^{\epsilon}\;\Prob\bigl[\calM(\calD')\in S\bigr] + \delta,
\label{eq:dp}
\end{equation}
а её композиция через моменты-бухгалтер позволяет точно отслеживать кумулятивный расход приватности $\epsilon$ по раундам обучения (используется в методах М2 и М7).

Модели пространства состояний переформулируют последовательную обработку графовых событий как линейную динамическую систему,
\begin{equation}
\bh_{t+1} = \bA\,\bh_t + \bB\,\bx_t, \qquad
\by_t = \bC\,\bh_t + \bD\,\bx_t,
\label{eq:ssm}
\end{equation}
чьё свёрточное ядро вычисляется за близкое к линейному время через быстрое преобразование Фурье (центральный элемент методов М4: MambaShield и М5: Стохастический трансформер).

Робастная оптимизация встраивает наихудшие возмущения графа непосредственно в оценку параметров через мин-макс целевую функцию
\begin{equation}
\min_\theta\;
\E_{(\bx,y)\sim\calD}
\Bigl[
  \max_{\|\bdelta\|\le\epsilon}
    \calL\!\bigl(f_\theta(\bx+\bdelta),\,y\bigr)
\Bigr],
\label{eq:robust_opt}
\end{equation}
чья структура двойственности связывает бюджет возмущений $\epsilon$ с чувствительностью выхода обученного отображения на графе (используется в методах М1 и М7).

Байесовский вывод формирует апостериорные предсказательные распределения, декомпозирующие полную предсказательную неопределённость на эпистемическую неопределённость (обусловленную ограниченностью обучающих данных) и алеаторную неопределённость (возникающую из-за присущего шума). В данной работе вариационное байесовское внимание и иерархические гауссовские процессы обеспечивают эту декомпозицию для графовых данных (центральный элемент методов М5 и М6).

Теоретико-игровое моделирование формализует состязательную обстановку как стратегическое взаимодействие между защитником (графовым классификатором) и атакующим. Анализ равновесия Штакельберга и алгоритмы онлайн-обучения без регрета определяют оптимальные стратегии защиты при модели наихудших возмущений (центральный элемент метода М7).

Византийско-устойчивая агрегация использует покоординатные усечённые средние для фильтрации повреждённых обновлений градиентов в условиях федеративного обучения на графах (центральный элемент метода М2: FedLLM-API).

В дополнение к этим математическим методам в исследовании применяются стандартные практики программной инженерии, проектирования распределённых систем и высокопроизводительных вычислений для обеспечения возможности обработки крупномасштабных темпоральных графовых потоков в рамках ограничений задержки операционных сред реального времени.


%%---------------------------------------------------------------------
\section*{Научная новизна}
%%---------------------------------------------------------------------

Научная новизна настоящей работы состоит из семи вкладов, каждый из которых сформулирован как принцип, концепция или теория, установленные автором.

\noindent 1.\quad \emph{Впервые} установлен принцип сертифицированной устойчивости моделей ГНС непрерывного времени, основанный на том, что динамика представлений узлов графа, управляемая дифференциальными уравнениями с ограниченной гладкостью, порождает доказуемые границы чувствительности выходных сигналов к возмущениям входных данных и структуры. Ни одна предшествующая работа не выводит сертифицированную робастность из гладкости управляющего дифференциального уравнения на темпорально эволюционирующих топологиях графов.

\noindent 2.\quad \emph{Оригинальная} теория совместного обеспечения византийской устойчивости, формальной приватности и распределительного выравнивания при обучении моделей ГНС на разрозненных фрагментах графов, принадлежащих разным организациям, без передачи исходных данных. Ни один предшествующий алгоритм не гарантирует совместно устойчивость к повреждениям, защиту приватности и выравнивание распределений для графовых данных.

\noindent 3.\quad \emph{Новая} концепция многоуровневого графового представления, в которой модели ГНС одновременно обучают структуру на нескольких семантических уровнях с механизмом защиты от утраты ранее усвоенных закономерностей при поступлении новых данных. Ни одна предшествующая архитектура не обучает многоуровневые темпоральные представления графов с одновременным предотвращением забывания.

\noindent 4.\quad \emph{Новый} подход к эффективному выводу моделей ГНС на основе селективных пространств состояний, позволяющий снизить вычислительную сложность с квадратичной до почти линейной при сохранении устойчивости к отравлению обучающих данных. Предшествующие структурированные модели пространства состояний не были адаптированы для графовых данных с формальными гарантиями робастности.

\noindent 5.\quad \emph{Впервые} предложена интегрированная теория калиброванной неопределённости для моделей ГНС, позволяющая разделять источники ошибок на устранимые (эпистемические) и неустранимые (алеаторные) и автоматически выявлять момент, когда модель перестаёт быть надёжной. Ни одна предшествующая работа не интегрирует эти два механизма неопределённости для графовой онлайн-адаптации.

\noindent 6.\quad \emph{Оригинальная} концепция единой сертификации устойчивости, в которой взаимодействие защитника и противника моделируется как стратегическая игра с доказуемо оптимальной стратегией выбора, связывающая все разработанные методы в целостную систему. Ни одна предшествующая единая система не обеспечивает доказуемые гарантии против всех четырёх типов возмущений одновременно.

\noindent 7.\quad \emph{Новая} теория предметно-ориентированного фундаментального моделирования на графовых структурах, обеспечивающая распознавание ранее невиданных категорий угроз без дополнительного обучения и совместимость с эволюционирующими протоколами. Эти вклады сведены к следующим основным научным результатам, выносимым на защиту.


%%---------------------------------------------------------------------
\section*{Основные научные результаты, выносимые на защиту}
%%---------------------------------------------------------------------

Результаты, выносимые на защиту, перечислены в порядке, отражающем их логические зависимости; каждый результат опирается на гарантии, установленные предшествующими.

(1) \textbf{Модели CT-TGNN и SDE-TGNN для сертифицированной динамики графов} --- оригинальные архитектуры непрерывного времени, в которых нейронные ОДУ и стохастические дифференциальные уравнения управляют эволюцией узлов и рёбер на динамических топологиях, обеспечивая сертифицированные границы робастности через анализ Липшица обученного векторного поля.

(2) \textbf{Алгоритм FedLLM-API для распределённого обучения графов с сохранением приватности} --- федеративный оптимизатор с доказанной сходимостью, сочетающий византийско-устойчивую агрегацию усечённым средним с $(\varepsilon{=}0.85,\,\delta{=}10^{-5})$ дифференциальной приватностью и выравниванием оптимального транспорта Вассерштейна для гетерогенных фрагментов графов.

(3) \textbf{Архитектура TripleE-TGNN для многоуровневого представления темпоральных графов} --- система тройного вложения, обучающая структуру на уровнях сервисов, трассировок и узлов с помощью двойного внимания и упругой консолидации весов, снижающая вычислительные затраты через структурированную разреженность и предотвращающая катастрофическое забывание.

(4) \textbf{Архитектура MambaShield для эффективного вывода на графах} --- селективная модель пространства состояний, снижающая сложность вывода с $\mathcal{O}(L^{2})$ до $\mathcal{O}(L\log L)$ через свёртки БПФ, с прогрессивной состязательной дистилляцией и PAC-байесовскими сертификатами обобщения.

(5) \textbf{Модели стохастического трансформера и UC-HGP для адаптации к дрейфу с учётом неопределённости} --- вариационное байесовское внимание с выводом на иерархических гауссовских процессах, обеспечивающее декомпозицию эпистемической и алеаторной неопределённости, онлайновое обнаружение дрейфа и гарантии сублинейного регрета при нестационарных графовых распределениях.

(6) \textbf{Единая система сертификации, интегрирующая модели (1)--(5)} --- единая методология оценки, связывающая все разработанные модели через теоретико-игровые равновесия Штакельберга и регуляризацию согласованности, валидированная на 6 бенчмарк-наборах данных общим объёмом 84,2 миллиона размеченных образцов.


%%---------------------------------------------------------------------
\section*{Теоретическая значимость}
%%---------------------------------------------------------------------

Теоретическая значимость работы состоит в создании единой математической основы для робастного обучения на динамических графовых структурах, интегрирующей пять областей математики, ранее не объединявшихся для задач на графах:

\noindent 1.\quad Нейронные ОДУ --- для сертифицированной динамики непрерывного времени на графах, при которой Липшицевость обученного отображения порождает доказуемые границы чувствительности.

\noindent 2.\quad Оптимальный транспорт --- для выравнивания распределений между организациями с одновременным обеспечением формальных гарантий дифференциальной приватности.

\noindent 3.\quad Селективные модели пространства состояний --- для снижения вычислительной сложности темпорального вывода с квадратичной до почти линейной, что делает обработку в реальном времени осуществимой.

\noindent 4.\quad Байесовский вывод --- для декомпозиции предсказательной неопределённости на устранимую (эпистемическую) и неустранимую (алеаторную) компоненты, позволяющей системе определять, когда её решения ненадёжны.

\noindent 5.\quad Теоретико-игровой выбор стратегий --- для моделирования взаимодействия защитника и противника как стратегической игры с доказуемо оптимальной стратегией защиты.

Эта основа устанавливает шесть формальных гарантий --- скорости сходимости ($\mathcal{O}(1/\sqrt{T})$), сертификаты робастности по Липшицу, границы $(\varepsilon,\delta)$-дифференциальной приватности, PAC-байесовские границы обобщения, сублинейный регрет ($R(T)\leq 2\sqrt{T\log|\mathcal{C}|}$) и единую сходимость (Теорема~6) --- которые являются предметно-независимыми и подтверждены как теоретически (через доказательства), так и эмпирически (через эксперименты на 84,2 миллионах образцов).


%%---------------------------------------------------------------------
\section*{Практическая значимость}
%%---------------------------------------------------------------------

Разработанные методы и алгоритмы применимы в нескольких областях и могут быть реализованы следующим образом.

\noindent 1.\quad \textbf{Здравоохранение.} Команды биомедицинского ИИ интегрируют CT-TGNN и SDE-TGNN в платформы поддержки клинических решений и лекарственного взаимодействия для обеспечения сертифицированно безопасных предсказаний на молекулярных графах и графах взаимодействия пациентов.

\noindent 2.\quad \textbf{Финансы.} Банки и финансовые регуляторы используют FedLLM-API для совместного обучения моделей обнаружения мошенничества на графах транзакций без обмена данными клиентов, сохраняя приватность через гарантии дифференциальной приватности.

\noindent 3.\quad \textbf{Промышленный IoT.} Инженеры OT-безопасности развёртывают MambaShield на периферийных шлюзах и контроллерах SCADA для обнаружения аномалий в реальном времени на графах датчиков благодаря его близкой к линейной вычислительной сложности.

\noindent 4.\quad \textbf{Автономные системы.} Инженеры безопасности используют UC-HGP для получения калиброванных эпистемических оценок уверенности для решений на графах слияния датчиков, что позволяет системе маркировать ненадёжные предсказания до того, как они приведут к опасным действиям.

\noindent 5.\quad \textbf{Кибербезопасность.} SOC-аналитики, инженеры кибербезопасности, исследователи MLSecOps и разработчики ИИ используют платформу с открытым исходным кодом RobustIDPS.ai (DOI: 10.5281/zenodo.19129512), интегрирующую 13+ нейросетевых моделей, 4-уровневую систему защиты LLM с обнаружением 94,5\% по 517 сценариям атак и работающую как плагин для Snort, Suricata и аналогичных систем СОПСВ.

Практическая применимость подтверждена актами внедрения от четырёх организаций: Государственной корпорации <<Росатом>> (защита ядерной инфраструктуры, точность 95,7\%), ООО <<Positive Technologies>> (федеративное обучение для коммерческих IDPS), Google Cloud Platform (безопасность на периферии, точность 96,8\%, снижение затрат на 40\%) и Suricata Open Source IDPS (адаптивное обнаружение с адаптацией к дрейфу).


%%---------------------------------------------------------------------
\section*{Личный вклад}
%%---------------------------------------------------------------------

Все результаты получены автором лично под научным руководством проф.\ А.\,Г.\,Трофимова. Автор лично разработал архитектуры CT-TGNN и SDE-TGNN, создал алгоритм FedLLM-API, спроектировал архитектуру TripleE-TGNN, разработал модель MambaShield, интегрировал стохастический трансформер с UC-HGP, построил единую систему сертификации и реализовал платформу RobustIDPS.ai. Вклад автора составляет 100\% в постановке задач, разработке методов, проведении экспериментов, создании программного обеспечения и подготовке публикаций.


%%---------------------------------------------------------------------
\section*{Достоверность научных результатов}
%%---------------------------------------------------------------------

Уверенность в представленных результатах обеспечивается несколькими независимыми линиями валидации.

Каждый предложенный алгоритм сопровождается формальным математическим анализом: теоремами сходимости с явными выражениями скорости, сертификатами робастности, выведенными из анализа чувствительности выхода графовой нейронной динамики и рандомизированного сглаживания, адаптированного к дискретным графовым структурам, гарантиями приватности, полученными через композицию моментов-бухгалтера гауссовского механизма, и границами вычислительной сложности, подтверждёнными как асимптотическим анализом, так и эмпирическим профилированием.

Эмпирическая оценка охватывает шесть бенчмарк-наборов данных, предоставляющих в совокупности более 84,2 миллиона размеченных образцов по 84 категориям возмущений. Используются три набора данных общего сетевого трафика: CIC-IoT-2023 (46,6 млн записей, 33 категории), CSE-CICIDS2018 (16,2 млн записей, 7 категорий) и UNSW-NB15 (2,5 млн записей, 9 категорий). Три специализированных набора данных облачных и периферийных сред из коллекции ICS3D дополняют их: Microsoft GUIDE (13,7 млн записей), Container Security (3,2 млн записей) и Edge-IIoT (2,0 млн записей).

Абляционные эксперименты систематически удаляют отдельные компоненты для количественной оценки маргинального вклада каждого. Состязательная робастность оценивается при шести семействах атак (FGSM, PGD, C\&W, DeepFool, гауссовский шум, отравление масок меток). Все эксперименты документированы таблицами гиперпараметров, спецификациями случайных затравок и описаниями оборудования, достаточными для независного воспроизведения.


%%---------------------------------------------------------------------
\section*{Апробация работы}
%%---------------------------------------------------------------------

Основные результаты настоящего исследования были представлены и прошли рецензирование на международных и национальных научных площадках:

\noindent 1.\quad Конференция IEEE молодых исследователей в области электротехники и электроники (ElCon-2024), ЭТУ-ЛЭТИ, Санкт-Петербург, IEEE Xplore.

\noindent 2.\quad Конференция IEEE молодых исследователей в области электротехники и электроники (ElCon-2025), ЭТУ-ЛЭТИ, Санкт-Петербург, IEEE Xplore.

\noindent 3.\quad Международная конференция NEUROINFORMATICS-2025, Springer.

\noindent 4.\quad III Всероссийская научно-техническая конференция <<Кибернетика и информационная безопасность>> (КИБ-2025), НИЯУ МИФИ, Москва.

\noindent 5.\quad 27-я Международная конференция IEEE молодых специалистов в области электронных устройств и материалов (EDM-2026), Республика Алтай.

\noindent 6.\quad 11-я конференция IEEE по технологиям интеграции электронных систем (IEEE ESTC 2026), Хельсинки, Финляндия.

\noindent 7.\quad Конференция IEEE <<Исследования и технологии для общества и промышленности>> (RTSI 2026), Эспоо, Финляндия.

\noindent 8.\quad NeurIPS 2026: 40-я ежегодная конференция по системам обработки нейронной информации, Сидней, Австралия.

Акты внедрения получены от четырёх организаций: Государственной корпорации <<Росатом>>, ООО <<Positive Technologies>>, Google Cloud Platform и Suricata Open Source IDPS.


%%---------------------------------------------------------------------
\section*{Публикации}
%%---------------------------------------------------------------------

Результаты настоящей диссертации изложены в рецензируемых публикациях в международных журналах и трудах конференций.

\subsection*{Опубликованные журнальные статьи}

\noindent
1. Anaedevha, R.\,N., Trofimov, A.\,G., Borodachev, Y.\,V. (2026). MambaShield: Selective State-Space Models for Poisoning-Resilient Sequential Modeling. \textit{Expert Systems with Applications}, Elsevier. DOI: 10.1016/j.eswa.2026.131175

\noindent
2. Anaedevha, R.\,N., Trofimov, A.\,G., Borodachev, Y.\,V. (2026). Uncertainty-calibrated hierarchical Gaussian processes for intrusion detection with multi-scale temporal modeling. \textit{Neurocomputing}, p.133105, Elsevier. DOI: 10.1016/j.neucom.2026.133105

\noindent
3. Anaedevha, R.\,N., Trofimov, A.\,G., Borodachev, Y.\,V. (2026). Temporal Adaptive Neural Ordinary Differential Equations with Deep Spatio-Temporal Point Processes for Real-Time Network Intrusion Detection. \textit{Complex and Intelligent Systems}, Springer. DOI: 10.1007/s40747-026-02291-7

\noindent
4. Anaedevha, R.\,N., Trofimov, A.\,G. (2024). Improved Robust Adversarial Model against Evasion Attacks on Intrusion Detection Systems. \textit{Optical Memory and Neural Networks (Information Optics)}, 33(Suppl.\,3), S414--S423. Springer. DOI: 10.3103/S1060992X24700681

\subsection*{Опубликованные труды конференций}

\noindent
5. Anaedevha, R.\,N., Trofimov, A.\,G. (2025). Stochastic Multimodal Transformer with Uncertainty Quantification for Robust Network Intrusion Detection. In: \textit{Advances in Neural Computation, Machine Learning, and Cognitive Research~IX. NEUROINFORMATICS~2025}. Studies in Computational Intelligence, vol.\,1241. Springer, Cham. DOI: 10.1007/978-3-032-07690-8\_35

\noindent
6. Anaedevha, R.\,N. \textit{et~al.}\ (2025). Integrated Deep Q-Networks and Actor-Critic Models for Enhanced Poisoning Attack Detection in Network Intrusion Detection Systems. In: \textit{2025 Conf.\ of Young Researchers in Electrical and Electronic Engineering (ElCon)}, ETU-LETI, Saint Petersburg. IEEE Xplore, Vol.\,33, pp.\,556--567.

\noindent
7. Anaedevha, R.\,N. \textit{et~al.}\ (2024). Application of Machine Learning Models for Device Identification in Wireless Network Traffic. In: \textit{2024 Conf.\ of Young Researchers in Electrical and Electronic Engineering (ElCon)}, Saint Petersburg. IEEE Xplore, pp.\,104--110. DOI: 10.1109/ElCon61730.2024.10468413

\subsection*{Принятые к публикации}

\noindent
8. Anaedevha, R.\,N., Trofimov, A.\,G., Borodachev, Y.\,V. (2025). Stochastic PAC-Bayesian Transformers for Network Intrusion Detection and Natural Language Processing Applications. \textit{Knowledge and Information Systems}, Springer. Принята.

\noindent
9. Anaedevha, R.\,N., Trofimov, A.\,G., Borodachev, Y.\,V. (2025). Byzantine-Resilient Stochastic Games for Federated Multi-Cloud Intrusion Detection. \textit{Complex and Intelligent Systems}, Springer. Принята.

\subsection*{Статьи на рецензировании}

\noindent
10. Anaedevha, R.\,N. \textit{et~al.}\ (2026). TripleE-TGNN: Triple-Embedding Temporal Graph Neural Networks for Multi-Granularity Microservices Security. \textit{IEEE Trans.\ Dependable and Secure Computing}.

\noindent
11. Anaedevha, R.\,N. \textit{et~al.}\ (2026). PQ-IDPS: Adversarially Robust Intrusion Detection for Post-Quantum Encrypted Traffic. \textit{IEEE Trans.\ Neural Networks and Learning Systems}.

\noindent
12. Anaedevha, R.\,N. \textit{et~al.}\ (2026). FedLLM-API: Privacy-Preserving Large Language Model Federation for Zero-Day API Threat Detection. \textit{IEEE Trans.\ Information Forensics and Security}.

\noindent
13. Anaedevha, R.\,N. \textit{et~al.}\ (2026). CT-TGNN: Continuous-Time Temporal Graph Neural Networks for Real-Time Threat Propagation Analysis. \textit{IEEE Trans.\ Services Computing}.

\subsection*{Сводка публикаций}

Общий объём составляет 18 статей: 4 опубликованы в журналах, 3 в трудах конференций, 2 приняты и находятся в печати, 9 на рецензировании. Шестнадцать статей (89\%) написаны в качестве первого автора. Публикации индексированы в Scopus, Web of Science и IEEE Xplore. Препринты доступны на TechRxiv и Research Square, записи цитирования поддерживаются в Google Scholar (h-индекс:~4) и ORCID (0000-0003-3649-1561).


%%---------------------------------------------------------------------
\section*{Соответствие паспорту специальности}
%%---------------------------------------------------------------------

Диссертация соответствует специальности 2.3.1 <<Системный анализ, управление и обработка информации, статистика>>:

\noindent П.4 <<Разработка методов и алгоритмов решения задач системного анализа, оптимизации, управления, принятия решений, обработки информации и искусственного интеллекта.>> --- Единая матрица $3\times4$ <<условия--требования>> декомпозирует проблему устойчивости динамических моделей ГНС на 7 подзадач; разработаны алгоритмы адъюнктного обучения CT-TGNN, FFT-свёртки MambaShield и византийски устойчивая оптимизация FedLLM-API; теоретико-игровые равновесия Штакельберга обеспечивают принятие решений с гарантиями; конвейер извлечения 83 признаков обрабатывает 84,2 млн образцов; 7 новых архитектур нейронных сетей с 6 доказанными теоремами.

\noindent П.5 <<Разработка специального математического и алгоритмического обеспечения систем анализа, оптимизации, управления, принятия решений, обработки информации и искусственного интеллекта.>> --- Сертификаты Липшица на основе неравенства Грёнуолла, пропагация моментов Фоккера--Планка, PAC-байесовские оценки и вассерштейновское OT-выравнивание составляют математическое обеспечение; 8 алгоритмов обеспечивают полный цикл от обучения до развёртывания; платформа RobustIDPS.ai (DOI: 10.5281/zenodo.19129512) с 46 страницами и 13+ моделями реализует программное обеспечение.


%%---------------------------------------------------------------------
\section*{Структура и объём работы}
%%---------------------------------------------------------------------

Диссертация состоит из Введения, пяти глав, Заключения, библиографического списка и трёх приложений. Основной текст составляет 210 страниц. Библиография содержит 190 наименований.

Во Введении обосновывается актуальность робастного и эффективного обучения на динамических графовых структурах как фундаментальной исследовательской задачи, представлены задокументированные свидетельства отказов моделей ГНС в четырёх критически важных областях, сформулирована цель исследования с семью последовательными задачами, определены объект и предмет исследования, описаны методы исследования с их формальными уравнениями, изложена научная новизна семи разработанных концепций и теорий, определены теоретическая и практическая значимость, представлен личный вклад автора, сформулированы шесть основных научных результатов, выносимых на защиту, подтверждена достоверность результатов, документированы мероприятия по апробации, публикации и соответствие паспорту специальности.

Глава~1, <<Анализ существующей литературы и выявление пробелов>>, представляет обзор литературы, следующий поэтапному охвату исследования: динамика непрерывного времени на эволюционирующих графах, распределённое федеративное обучение с приватностью, вычислительная эффективность для крупномасштабных графов и нестационарная адаптация с оценкой неопределённости. Каждый раздел завершается связующим утверждением, соединяющим выявленный пробел с решением, разработанным в следующей главе. Выявлено семь исследовательских пробелов из 128 цитируемых работ.

Глава~2, <<Методы, алгоритмы и математические основы>>, представляет полные математические выводы, доказательства сходимости и робастности, алгоритмический псевдокод, архитектурные схемы и анализ сложности для всех семи методов (М1--М7) и их расширений (SDE-TGNN, CyberSecLLM). Единая постановка задачи последовательно выводится в 7 этапах и декомпозируется на 7 подзадач. 62 уравнения, 6 теорем, 8 алгоритмов.

Глава~3, <<Адаптация и согласование с гибридными системами обнаружения и предотвращения вторжений>>, является предметно-ориентированной главой, в которой математические методы Главы~2 отображаются на сущности и ограничения кибербезопасности. Определена таксономия из 34 классов, описан конвейер 83 инженерных признаков, представлены гиперпараметры и обоснование их выбора для каждой модели.

Глава~4, <<Экспериментальные результаты и валидация>>, представляет полную экспериментальную программу на 6 наборах данных (84,2 млн образцов, 84 категории, 23 метрики), включая результаты по отдельным методам, абляцию, состязательную робастность при 6 семействах атак, кросс-доменный перенос и федеративное обучение при византийских повреждениях.

Глава~5, <<RobustIDPS.ai: Программная платформа для промышленного развёртывания>>, описывает платформу с открытым исходным кодом (DOI: 10.5281/zenodo.19129512): 46 страниц, 9 функциональных групп, 13+ моделей, 4-уровневая защита LLM (94,5\%), модуль Multi-Agent PQC-IDS и интеграция со Snort/Suricata. Пользовательское исследование подтвердило 34\% снижение времени сортировки оповещений.

В Заключении перечислены все научные результаты, подтверждено покрытие всех 12 ячеек матрицы условие--требование, указаны акты внедрения и намечены направления дальнейших исследований.

Приложение~А содержит полные доказательства всех теорем. Приложение~Б предоставляет таблицы гиперпараметров и спецификации инфраструктуры. Приложение~В документирует дополнительные экспериментальные результаты.
